\section{Conclusion}
\label{sec:conclusion}

A constraint inferred from data describes what the data happens to satisfy, whereas a constraint contained in governing text describes what the data is required to satisfy. This distinction \subs{ }{}has received limited attention in database research, even though many published datasets are accompanied by documents that define data quality criteria at some level. 

In this paper, we show that the gap between data and normative text can be addressed automatically. Our results indicate that lightweight classification is sufficient to detect \subs{ }{}most checkable requirements, and that a structured multi-step pipeline grounded in both text and schema can translate them into executable constraints while avoiding hallucinated predicates and literals. 

This enables a new form of validation: \system\ can reveal compliance failures even when they do not manifest as discoverable \subs{pattern}{patterns} in the data. In such settings, the normative document becomes \subs{ }{}part of the data specification, and keeping that specification aligned with evolving text emerges \add{as} a concrete challenge for data management. \add{Because each constraint comes from the text, it remains valid while the regulation is unchanged and can be re-executed, also incrementally~\cite{youri2024}, as the data evolves; only a change in the text requires a new mapping. Our datasets have at most 158 sentences, whereas legal documents often have hundreds, with nested clauses that must first be segmented and resolved; since cost grows linearly with the number of sentences, we expect the pipeline to scale to them, but this remains to be shown.}


